\originalchapter{二能级共振与相干控制}\label{ch:res}
\originalsection{有效哈密顿量与旋转参考系}
原子有许多能级, 但窄带外场只接近一个跃迁时, 可把其他能级的主要作用归入小的能级修正. 令$\ket e,\ket g$相隔$\hbar\omega_0$, 定义$\sigma_z=\ket e\bra e-\ket g\bra g$, $\sigma_+=\ket e\bra g$. 外场的偶极矩阵元规定为$-\bra e\mathbf d\cdot\mathbf E_0\ket g=\hbar\Omega$, 取相位使$\Omega$实. 实场$\mathbf E_0\cos\omega_Lt$因余弦分解而产生$\hbar\Omega/2$的共振耦合. 本卷$\Omega$是角频率, 普通频率为$\Omega/(2\pi)$.

在$U=\exp(-\ii\omega_Lt\sigma_z/2)$定义的参考系中, $\ket\psi=U\ket{\widetilde\psi}$满足$\widetilde H=U^\dagger HU-\ii\hbar U^\dagger\dot U$. 保留慢项得
\begin{equation}\widetilde H=\frac{\hbar}{2}\begin{pmatrix}-\Delta&\Omega\\\Omega&\Delta\end{pmatrix},\qquad\Delta=\omega_L-\omega_0.\label{eq:rabiH}\end{equation}
省去的项以$2\omega_L$振荡. 旋波近似要求$|\Delta|,|\Omega|\ll\omega_0$且关注时间远大于一个光学周期; 它不是对任意强场或极短脉冲都成立的恒等变换. 取$\Omega_R=\sqrt{\Delta^2+\Omega^2}$, 因$\widetilde H^2=\hbar^2\Omega_R^2I/4$, 指数级数的奇偶项分别相加, 得
\begin{equation}\widetilde U(t)=I\cos\frac{\Omega_Rt}{2}-\frac{2\ii\widetilde H}{\hbar\Omega_R}\sin\frac{\Omega_Rt}{2},\qquad P_e(t)=\frac{\Omega^2}{\Omega_R^2}\sin^2\frac{\Omega_Rt}{2}.\label{eq:rabi}\end{equation}
最后一式假定初态为$\ket g$. 共振时面积$\int\Omega(t)\dd t=\pi$的脉冲翻转布居; $\pi/2$脉冲制备等幅叠加. 失谐既加快相位演化, 也缩小可达的激发布居. 弱驱动的短时概率为$\Omega^2t^2/4$, 不能直接用线性速率方程替代.

\originalsection{Bloch向量、Ramsey与回波}
二能级密度矩阵可写成$\rho=(I+\mathbf r\cdot\boldsymbol\sigma)/2$, 其中$r_i=\Tr(\rho\sigma_i)$是实数, $|\mathbf r|\le1$. 把式\eqref{eq:rabiH}写成$\hbar\mathbf b\cdot\boldsymbol\sigma/2$, 利用$[\sigma_i,\sigma_j]=2\ii\epsilon_{ijk}\sigma_k$得$\dot{\mathbf r}=\mathbf b\times\mathbf r$. 纯态在球面, 混态在球内; 投影到$z$轴给布居差, 横向投影给相干性.

两个同相$\pi/2$脉冲之间自由等待$T$, 在旋转系中产生相对相位$\Delta T$. 直接把两个脉冲矩阵与自由演化矩阵相乘, 可得$P_e=[1+\cos(\Delta T)]/2$; 第二脉冲相位改变$\phi$时, 余弦的自变量变为$\Delta T-\phi$. 窄条纹由自由等待时间决定, 而脉冲的有限带宽决定外包络. 因此频率标准可以用短脉冲控制, 用长等待提高频率辨别力.

静态但位置不同的失谐使横向向量散开. 在$T/2$插入$\pi$脉冲, 第一段的快慢相位顺序翻转, 第二段抵消第一段. 用$\sigma_x\sigma_z\sigma_x=-\sigma_z$即可验证总自由相位为零. 回波能撤销可逆的非均匀退相干, 不能追回已泄露到环境且未被逆转的信息. 这一区别在第\ref{ch:open}章的$T_1,T_2$中定量出现.

\originalsection{绝热跟随与能级反交叉}
矩阵的本征能量为$\pm\hbar\Omega_R/2$. 可用角度$\tan\theta=\Omega/(-\Delta)$描述本征轴. 随着失谐缓慢扫过零, 一个瞬时本征态的主要布居由$g$转移到$e$. 在瞬时本征基中, 时间变化还产生非对角项$\hbar\dot\theta/2$, 所以局域绝热条件是$|\dot\theta|\ll\Omega_R$. 有限耦合使最小能隙$\hbar|\Omega|$不为零; 扫过太快便产生跃迁.

对线性扫频$\Delta=vt$, 从$t=-\infty$开始的Landau--Zener穿越给出绝热失败概率$\exp[-\pi\Omega^2/(2|v|)]$. 此精确特殊函数结果在本章作为已注明模型的工具使用, 不把局域绝热估计当作该指数的证明. 相比精确面积的$\pi$脉冲, 扫频转移对部分幅度误差更稳健, 代价是时间更长及对退相干更敏感.
